\begin{figure}[htbp]
\centering
\begin{tikzpicture}[causal,x=1cm,y=1cm,every node/.style={font=\small,inner sep=2pt}]
\node (c06x) at (0,3) {$X$};
\node (c06b) at (0,2) {$B$};
\node (c06d) at (0,1) {$D$};
\node (c06y) at (0,0) {$Y$};
\draw[->,blue!65!black] (c06x) -- node[right] {$b$} (c06b);
\draw[->,blue!65!black] (c06b) -- node[right] {$d$} (c06d);
\draw[->,blue!65!black] (c06d) -- node[right] {$h$} (c06y);
\end{tikzpicture}
\origfig{6}{间接效应。}
\label{fig:ch01c-06}
\end{figure}
